\section{\texorpdfstring{A static curve in \(r/\ell_W\)}{A static curve in r/ellW}}
\label{sec:h-dimevol}

On a pure power the logarithmic derivative is the constant \(d\),
independent of \(\ell\).
On the chain it is a function of the dimensionless ratio \(r/\ell_W\):
\begin{equation}
\label{eq:deff}
d_{\mathrm{eff}}(r)\to
\begin{cases}
2 & r/\ell_W\to 0,\\
3 & 1\ll r/\ell_W\ll S,\\
2 & r/\ell_W\gg S.
\end{cases}
\end{equation}
An observer at fixed graph distance sees a fixed exponent.
An observer who changes the probe scale, on a fixed graph, walks along
\eqref{eq:deff}.
That is a scale dependence.
A time dependence at fixed probe radius needs \(\ell_W(t)\) or \(S(t)\),
hence a changing graph (Section~\ref{sec:h-de}).

Until the graph changes, the small-\(r\) and large-\(r\) ends of
\eqref{eq:deff} are two ends of one static curve.
You could collapse \(d_{\mathrm{eff}}\) against \(\log(r/\ell_W)\) for
several \(\lambda\).
A collapse onto one curve means the single length \eqref{eq:ell} controls
the crossover.
A failure to collapse means that formula does not.
Section~\ref{sec:h-cross} is the same length read from the return
probability rather than from the ball.
